\section{Replication, Cluster Formation y Hierarchical Sparsity}

GLOM replica el mismo embedding de alto nivel en las muchas posiciones que cubre un object, lo que parece un gran desperdicio. La refutación de Hinton es esta: mientras el binding no está determinado, cada posición necesita conservar una hypothesis independiente para decidir de forma gradual «qué posiciones deben ser iguales»; la replicación no es la compresión óptima final, sino el costo de la inference flexibility y de la locality.

\lecturefigure{slide-30-object-embedding-replication.jpg}{Cada object patch replica el object-level embedding; parece un desperdicio, pero conserva la hypothesis local.}{Video oficial de Stanford Online 00:48:52--00:49:47.}

\teachervoice{El ponente recurre a una analogía de biology: las células replican el mismo ADN, y distintas partes de un órgano pueden tener protein-expression vectors similares. La función didáctica de esta analogía es mostrar que la locality puede justificar la redundancia; no afirma que un embedding de GLOM equivalga a un gen o a la expresión de proteínas.}

\subsection{Forming clusters, no clustering sobre puntos fijos}

El clustering ordinario recibe data points fijos y después descubre los grupos; en GLOM, los propios data points cambian con la interacción bottom-up, top-down, lateral y temporal. El sistema no agrupa embeddings ya existentes, sino que forma conjuntamente los embeddings y las cluster boundaries dentro de la recurrent dynamics.

\begin{importantbox}{El significado computacional de «Hedge your bets»}
En los primeros steps, la posición $x$ puede conservar una parent hypothesis distinta de la de sus vecinas; solo a medida que se acumula evidencia, las positions compatibles se absorben en la misma island. Si desde el inicio se comprimiera todo el object en una sola shared variable, un binding erróneo sería difícil de corregir localmente. Lo que se obtiene a cambio de la replicación es delayed commitment.
\end{importantbox}

\begin{warningbox}{La formación dinámica no implica obtener automáticamente los discrete clusters correctos}
Los vectores continuos pueden formar transiciones difusas, varios attractors locales o un global consensus. Si se necesitan object slots discretos, hay que definir además el readout, el threshold, la connectedness, el temporal tracking y la birth/death rule. Esta clase no completa ese contract.
\end{warningbox}

\subsection{Los niveles altos, más lejanos y más dispersos}

En los niveles bajos, las parts son pequeñas y sus bordes son finos, por lo que requieren dense interaction de corto alcance; en los niveles altos, las islands de object/scene son grandes y su redundancia interna es alta, de modo que pueden ampliar el receptive radius y muestrear solo unas cuantas posiciones. El objetivo es que, al subir de hierarchy level, el communication range crezca sin que se dispare el edge count.

\lecturefigure{slide-31-sparse-high-level-replication.jpg}{Las islands de alto nivel son más grandes y pueden usar attention connections de mayor alcance y más dispersas.}{Video oficial de Stanford Online 00:49:48--00:50:07.}

\begin{knowledgebox}{Lectura de la figura: que el range aumente no exige que el edge count crezca en la misma proporción}
Una object island de alto nivel cubre muchos patches y se espera que sus vectores internos sean muy redundantes; por eso, una posición no necesita conectarse con todas las posiciones de la island: basta con que un sparse sample alcance puntos representativos para que pueda obtener la misma whole hypothesis. La figura expresa una range--density tradeoff: el interaction radius crece con el level y la densidad de conexiones disminuye. No proporciona el neighbor-search cost, la miss rate ni el throughput real, así que solo puede tomarse como una systems hypothesis.
\end{knowledgebox}

\begin{table}[H]
\centering
\small
\caption{Intención de diseño de la locality, el range y la sparsity en la jerarquía de GLOM.}
\begin{tabularx}{\textwidth}{L{0.17\textwidth}L{0.24\textwidth}L{0.24\textwidth}>{\raggedright\arraybackslash}X}
\toprule
Nivel & Entidad típica & Interaction pattern & Riesgo principal \\
\midrule
Bajo & edge, texture, small part & De corto alcance, relativamente dense; protege los bordes finos & Ver solo lo local puede hacer perder el whole context \\
Medio & major part, object component & De alcance medio, gated por similarity & Tiende a fusionar instancias parecidas pero distintas \\
Alto & object, person, scene & De largo alcance, relativamente sparse; aprovecha la redundancia de las islands grandes & El sparse sample puede omitir objetos pequeños o generar conflictos entre varias instancias \\
\bottomrule
\end{tabularx}
\end{table}

\begin{warningbox}{«La misma carga de trabajo en cada nivel» es un objetivo, no un resultado medido}
Para verificar esta afirmación hay que precisar, en cada nivel, el número de posiciones, el neighbor count, el embedding width, el número de pasos de iteración, la routing/ANN search, el memory traffic y el communication pattern. Decir solo «más sparse» no basta para demostrar que los FLOPs, la latency o la energy se mantienen constantes.
\end{warningbox}

\subsection{Resumen de la sección}

La embedding replication conserva grados de libertad para el gradual binding y la local correction; la hierarchical sparsity intenta convertir esa redundancia en una interacción escalable. Juntas expresan la disyuntiva central de GLOM: es preferible repetir en la state memory antes que comprimir de antemano una uncertain structure en un único node irreversible.
